% p140_1; source PDF page 140; manually verified editable table.
\begin{tikzpicture}[x=1cm,y=-1cm,inner sep=0pt,outer sep=0pt]
\path[use as bounding box] (0,0) rectangle (16.000,10.975);
\draw[line width=.4pt] (0.000,0.000) rectangle (5.500,0.855);
\node[align=center,text width=5.240cm,font={\linespread{1}\fontsize{10}{13.50}\selectfont\bfseries}] at (2.750,0.427) {案情};
\draw[line width=.4pt] (5.500,0.000) rectangle (7.500,0.855);
\node[align=center,text width=1.740cm,font={\linespread{1}\fontsize{10}{13.50}\selectfont\bfseries}] at (6.500,0.427) {查证情况};
\draw[line width=.4pt] (7.500,0.000) rectangle (16.000,0.855);
\node[align=center,text width=8.240cm,font={\linespread{1}\fontsize{10}{13.50}\selectfont\bfseries}] at (11.750,0.427) {结论};
\draw[line width=.4pt] (0.000,0.855) rectangle (5.500,2.860);
\node[align=left,text width=5.240cm,font={\linespread{1}\fontsize{10}{13.50}\selectfont}] at (2.750,1.857) {（1）故意的同时犯。甲、乙没有意思联络，故意杀人，同时向丙开枪，丙死亡};
\draw[line width=.4pt] (0.000,2.860) rectangle (5.500,4.865);
\node[align=left,text width=5.240cm,font={\linespread{1}\fontsize{10}{13.50}\selectfont}] at (2.750,3.862) {（2）过失的同时犯。甲、乙没有意思联络，各自过失导致子弹飞向丙，丙死亡};
\draw[line width=.4pt] (0.000,4.865) rectangle (5.500,6.870);
\node[align=left,text width=5.240cm,font={\linespread{1}\fontsize{10}{13.50}\selectfont}] at (2.750,5.867) {（3）故意的共同正犯。甲、乙有意思联络，共同故意，同时向丙开枪，丙死亡};
\draw[line width=.4pt] (0.000,6.870) rectangle (5.500,10.975);
\node[align=left,text width=5.240cm,font={\linespread{1}\fontsize{10}{13.50}\selectfont}] at (2.750,8.923) {（4）共同过失案：甲、乙有意思联络，共同过失，同时向丙开枪，丙死亡（打野猪案）};
\draw[line width=.4pt] (5.500,0.855) rectangle (7.500,10.975);
\node[align=left,text width=1.740cm,font={\linespread{1}\fontsize{10}{13.50}\selectfont}] at (6.500,5.915) {事后查明，只有一枪致命，但无法查明致命一枪是谁打的};
\draw[line width=.4pt] (7.500,0.855) rectangle (16.000,2.860);
\node[align=left,text width=8.240cm,font={\linespread{1}\fontsize{10}{13.50}\selectfont}] at (11.750,1.857) {甲、乙不构成共同正犯，需单独处理。根据存疑时有利于被告原则，甲、乙都构成故意杀人罪未遂，而非既遂};
\draw[line width=.4pt] (7.500,2.860) rectangle (16.000,4.865);
\node[align=left,text width=8.240cm,font={\linespread{1}\fontsize{10}{13.50}\selectfont}] at (11.750,3.862) {甲、乙不构成共同正犯，需单独处理。根据存疑时有利于被告原则，甲、乙都不构成过失致人死亡罪，都无罪};
\draw[line width=.4pt] (7.500,4.865) rectangle (16.000,6.870);
\node[align=left,text width=8.240cm,font={\linespread{1}\fontsize{10}{13.50}\selectfont}] at (11.750,5.867) {由于甲、乙是共同正犯，根据“部分实行、全部负责”原则，甲、乙对死亡结果都要负责，都构成故意杀人罪既遂};
\draw[line width=.4pt] (7.500,6.870) rectangle (16.000,10.975);
\node[align=left,text width=8.240cm,font={\linespread{1}\fontsize{10}{13.50}\selectfont}] at (11.750,8.923) {（1）根据要求故意且无罪说（少数说），甲、乙不构成共同正犯，基于“存疑时有利于被告”原则，都不构成过失致人死亡罪，都无罪。（2）根据要求故意且有罪说（多数说），甲、乙相互给对方的致人死亡的结果都有心理性贡献，因此对对方的致人死亡应负责。甲乙不构成共同正犯，各自构成过失致人死亡罪};
\end{tikzpicture}
